\documentclass[reqno]{exam}
\documentclass[reqno]{exam}

\usepackage{amssymb, amsmath, amsthm, stmaryrd}
\usepackage{fullpage, parskip}
\usepackage{graphicx}
\usepackage{enumitem}

% Improved spacing
\usepackage[top=1in, bottom=0.8in, left=1in, right=1in]{geometry}


\title{\normalsize{
  JHED: \underline{\hspace{4cm}} \hspace{1em} Name: \underline{\hspace{9cm}}}
   }
\author{}
\date{}

% Custom spacing for better page utilization
\renewcommand{\baselinestretch}{0.95}

\begin{document}
\maketitle
\vspace{-2cm} % Reduce space after title

\begin{questions}
  \question[8]
Suppose you have a machine that works in \textit{base 10} instead of \textit{base 2} with precision $d = 1$ (number of digits in the fractional part). Let $\mathbb{F}$ denote the set of all floating point numbers representable on this machine.

\begin{parts}
  \part What is the machine epsilon $\epsilon_{\text{mach}}$?
  \begin{solution}
  Machine epsilon is defined as $\epsilon_{\text{mach}} = \beta^{-d}$ where $\beta$ is the base and $d$ is the precision (number of digits after the leading digit).
  
  With $\beta = 10$ and $d = 1$:
  $$\epsilon_{\text{mach}} = 10^{-1} = \boxed{0.1}$$
  \end{solution}
  \vfill
  \part How many elements of $\mathbb{F}$ are there in the real interval $[50, 5000]$, including the endpoints?
  \begin{solution}
  With base 10 and precision $d=1$, elements of $\mathbb{F}$ have the form $a.b \times 10^n$ where $a \in \{1,...,9\}$ and $b \in \{0,1,...,9\}$.
  
  \textbf{In $[50, 99]$:} $50, 51, ..., 99$ $\rightarrow$ 50 elements
  
  \textbf{In $[100, 990]$:} $100, 110, 120, ..., 990$ $\rightarrow$ 90 elements
  
  \textbf{In $[1000, 5000]$:} $1000, 1100, 1200, ..., 5000$ $\rightarrow$ 41 elements
  
  Total: $50 + 90 + 41 = \boxed{181}$ elements
  \end{solution}
  \vfill
  \part What is the element of $\mathbb{F}$ closest to the real number $\pi$? 
  \begin{solution}
    $\boxed{3.1}$.
  \end{solution}
  \vfill
  \part What is a stretch of 10 consecutive positive integers not in $\mathbb{F}$? (more than one answer possible)
  \begin{solution}
   One possible answer: $\boxed{1001, 1002, 1003, 1004, 1005, 1006, 1007, 1008, 1009, 1010}$
  \end{solution}
  \vfill
\end{parts}




  \newpage
   Let $a = 5.65$, $b = 6.78$, $c = 2.34$.

    \question[6]


    Suppose $\mathrm{fl}(x)$ is the number in $\mathbb{F}$ \textbf{closest} to the real number $x$ (round half up).
Compute the following, showing intermediate steps:

\begin{parts}
  \part $a + (b - c)$
  \begin{solution}
  \textbf{Step 0:} Convert to $\mathbb{F}$: $\mathrm{fl}(a) = 5.7$, $\mathrm{fl}(b) = 6.8$, $\mathrm{fl}(c) = 2.3$
  
  \textbf{Step 1:} Compute $b - c = 6.8 - 2.3 = 4.5$ (already in $\mathbb{F}$)
  
  \textbf{Step 2:} Compute $a + 4.5 = 5.7 + 4.5 = 10.2$
  
  \textbf{Step 3:} Round to $\mathbb{F}$: $\mathrm{fl}(10.2) = 10$ (since elements in this range are $10, 11, 12, ...$)
  
  $$\boxed{10}$$
  \end{solution}
  \vfill
  \part $(a + b) - c$
  \begin{solution}
  \textbf{Step 0:} Convert to $\mathbb{F}$: $\mathrm{fl}(a) = 5.7$, $\mathrm{fl}(b) = 6.8$, $\mathrm{fl}(c) = 2.3$
  
  \textbf{Step 1:} Compute $a + b = 5.7 + 6.8 = 12.5$
  
  \textbf{Step 2:} Round to $\mathbb{F}$: $\mathrm{fl}(12.5) = 13$ (round half up)
  
  \textbf{Step 3:} Compute $13 - c = 13 - 2.3 = 10.7$
  
  \textbf{Step 4:} Round to $\mathbb{F}$: $\mathrm{fl}(10.7) = 11$
  
  $$\boxed{11}$$
  \end{solution}
  \vfill
\end{parts}

    \question[6]

    Suppose instead $\mathrm{fl}(x)$ is the number in $\mathbb{F}$ obtained by truncating the real number $x$ (i.e., chopping off digits beyond the precision).
Compute the following, showing intermediate steps:

\begin{parts}
  \part $a + (b - c)$
  \begin{solution}
  \textbf{Step 0:} Convert to $\mathbb{F}$: $\mathrm{fl}(a) = 5.6$, $\mathrm{fl}(b) = 6.7$, $\mathrm{fl}(c) = 2.3$
  
  \textbf{Step 1:} Compute $b - c = 6.7 - 2.3 = 4.4$ (already in $\mathbb{F}$)
  
  \textbf{Step 2:} Compute $a + 4.4 = 5.6 + 4.4 = 10.0$ (already in $\mathbb{F}$)
  
  $$\boxed{10}$$
  \end{solution}
  \vfill
  \part $(a + b) - c$
  \begin{solution}
  \textbf{Step 0:} Convert to $\mathbb{F}$: $\mathrm{fl}(a) = 5.6$, $\mathrm{fl}(b) = 6.7$, $\mathrm{fl}(c) = 2.3$
  
  \textbf{Step 1:} Compute $a + b = 5.6 + 6.7 = 12.3$
  
  \textbf{Step 2:} Truncate to $\mathbb{F}$: $\mathrm{fl}(12.3) = 12$
  
  \textbf{Step 3:} Compute $12 - c = 12 - 2.3 = 9.7$ (already in $\mathbb{F}$)
  
  $$\boxed{9.7}$$
  \end{solution}
  \vfill
\end{parts}


\end{questions}
\end{document}